\documentclass[11pt]{article}
\usepackage[a4paper,margin=18mm]{geometry}
\usepackage{fontspec}
\setmainfont{Latin Modern Roman}
\usepackage{amsmath,amssymb,amsthm,amscd,graphicx,color}
\usepackage{xurl}
\usepackage[unicode,hidelinks]{hyperref}
\allowdisplaybreaks
\setlength\emergencystretch{3em}

\makeatletter
\newtheorem{theorem}{Theorem}
\newtheorem*{theorem*}{Theorem}
\newtheorem{lemma}{Lemma}
\newtheorem*{lemma*}{Lemma}
\newtheorem{corollary}{Corollary}
\newtheorem*{corollary*}{Corollary}
\newtheorem{proposition}{Proposition}
\newtheorem*{proposition*}{Proposition}
\newtheorem{conjecture}{Conjecture}
\newtheorem*{conjecture*}{Conjecture}
{\theoremstyle{definition} \newtheorem{definition}{Definition}
\newtheorem*{definition*}{Definition}
\newtheorem{example}{Example}
\newtheorem*{example*}{Example}
\newtheorem{remark}{Remark}
\newtheorem*{remark*}{Remark}
\newtheorem{note}{Note}
\newtheorem*{note*}{Note}
}

\makeatother
\DeclareMathOperator{\ev}{ev}
\DeclareMathOperator{\gr}{gr}
\DeclareMathOperator{\spn}{span}

\numberwithin{equation}{section}


\begin{document}
\begin{center}\Large Exact Demazure quotient sequence and mixed level-one-Demazure/adjoint fusion identification for a finite-dimensional complex simple Lie algebra, k>=\allowbreak{}1 and 0<=\allowbreak{}i<=\allowbreak{}k\end{center}
\noindent\textbf{Stable ID:} 2027\par
\noindent\textbf{Authors:} Bhimarthi Ravinder\par
\noindent\textbf{Original work:} Demazure Modules, Chari-Venkatesh Modules and Fusion Products\par
\noindent\textbf{Version:} Official SIGMA 10 (2014) article 110, DOI 10.3842/SIGMA.2014.110; journal PDF and native TeX acquired 2026-10-01, exact bytes pinned by SHA256\par
\noindent\textbf{Selected task scope:} Full Theorem1 parts1-2 as a single main proof unit, with all displayed quotient exponents and shift unchanged; mixed fusion has k+\allowbreak{}1-i level-one D(1,theta) factors and i evaluation adjoint factors. No simply-laced restriction imposed; no extra CV or truncated-Weyl consequence selected.\par
\noindent\textbf{Difficulty:} H2: The proof constructs the kernel map from a highest-root current vector and proves its defining relations through nontrivial root-string and Garland current identities. A fusion-product surjection and a two-parameter dimension induction then force both exactness and the mixed fusion identification for every Lie type, rather than just specializing a same-level fusion theorem.\par
\noindent\textbf{Editorial boundary note (not author text):} The complete original two-clause Theorem 1 is supplied as one exact-sequence and mixed-fusion main problem unit for every finite-dimensional complex simple Lie algebra, k>=\allowbreak{}1 and 0<=\allowbreak{}i<=\allowbreak{}k. The exact grading shift and single highest-root-current quotient generators are retained. All new local presentation, quotient-kernel, mixed-fusion relation and dimension-bound arguments and both induction directions are included, with arbitrary pairwise distinct fusion parameters. Garland and the all-adjoint fusion theorem are explicitly published prior prerequisites. The source Cartan relation slip and the later authoritative local presentation remain unchanged with their notice. No CV restatement, truncated-Weyl corollary, simply-laced restriction or added parameter/type variant is selected.\par
\noindent\textbf{Source:} \url{https://sigma-journal.com/2014/110/}\quad DOI: \url{https://doi.org/10.3842/SIGMA.2014.110}\par
\noindent\textbf{License and attribution:} Copyright retained by the named authors. Published by SIGMA under CC BY 4.0: \url{https://creativecommons.org/licenses/by/4.0/}. Source: \url{https://sigma-journal.com/about.html}. Adaptation consists of standalone excerpt formatting, cover and numbering restoration; original author language and mathematical text are retained.\par
\noindent\textbf{Editorial symbol notice (not author text):} Definition1 has an extra trailing=\allowbreak{}0 after its Cartan-weight relation, while Proposition1 and(4.2) explicitly give the correct same weight relation used in the proof. Original definition and authoritative local presentation remain unchanged.\par
\medskip\hrule\medskip\noindent\textbf{Author text begins below.}\par
\setcounter{section}{1}
\setcounter{subsection}{0}
\setcounter{subsubsection}{0}
\setcounter{equation}{0}
% Original source lines 90--110
\begin{theorem}
\label{MT}
Let $k\geq 1$.
For $0\leq i \leq k$, we have the following:
\begin{enumerate}\itemsep=0pt
\item[$1)$] a~short exact sequence of $\mathfrak{g}[t]$-modules,
\begin{gather*}
0\rightarrow \tau_{2k+1-i} \big(D(1, k\theta)/\big\langle \big(x^{-}_\theta \otimes t^{2k-i}\big) \overline{w}_{k\theta}\big\rangle\big)
\\
\phantom{0}
\xrightarrow{\phi^{-}} D\big(1,(k+1)\theta\big)/\big\langle \big(x^{-}_\theta \otimes t^{2k+2-i}\big) \overline{w}_{(k+1)\theta}\big\rangle
\\
\phantom{0}
\xrightarrow{\phi^+} D\big(1,(k+1)\theta\big)/\big\langle \big(x^{-}_\theta \otimes t^{2k+1-i}\big) \overline{w}_{(k+1)\theta}\big\rangle \rightarrow 0;
\end{gather*}
\item[$2)$] an isomorphism of $\mathfrak{g}[t]$-modules,
\begin{gather*}
D\big(1,(k+1)\theta\big)/\langle \big(x^{-}_\theta \otimes t^{2k+2-i}\big) \overline{w}_{(k+1)\theta}\rangle \cong D(1,\theta)^{* (k+1-i)}* \textup{ev}_0 V(\theta)^{*i}.
\end{gather*}
\end{enumerate}
\end{theorem}
\setcounter{section}{1}
\setcounter{subsection}{0}
\setcounter{subsubsection}{0}
\setcounter{equation}{0}
% Original source lines 168--731
\section{Preliminaries}\label{section2}

Throughout the paper, $\mathbb{C}$ denote the f\/ield of complex numbers, $\mathbb{Z}$ the set of integers,
$\mathbb{Z}_{\geq 0}$ the set of non-negative integers, $\mathbb{N}$ the set of positive integers and $\mathbb{C}[t]$
the polynomial ring in an indeterminate~$t$.

{\bf 2.1.}~Let $\mathfrak{a}$ be a~complex Lie algebra, $\mathbf{U}(\mathfrak{a})$ the corresponding universal
enveloping algebra.
The current algebra associated to $\mathfrak{a}$ is denoted by $\mathfrak{a}[t]$ and def\/ined as $\mathfrak{a} \otimes
\mathbb{C}[t]$, with the Lie bracket
\begin{gather*}
[a \otimes t^r, b \otimes t^s]=[a, b]\otimes t^{r+s},
\qquad
\text{for all}
\quad
a, b \in \mathfrak{a}
\quad
\text{and}
\quad
r, s \in \mathbb{Z}_{\geq 0}.
\end{gather*}
We let $\mathfrak{a}[t]_{+}$ be the ideal $\mathfrak{a}\otimes t\mathbb{C}[t]$. The degree grading on $\mathbb{C}[t]$
gives a~natural $\mathbb{Z}_{\geq 0}$-grading on~$\mathbf{U}(\mathfrak{a}[t])$ and the subspace of grade~$s$ is given~by
\begin{gather*}
\mathbf{U}(\mathfrak{a}[t])[s]= \spn \Big\{(a_1\otimes t^{r_1})\cdots (a_k\otimes t^{r_k}):k\geq 1,\,
a_i\in\mathfrak{a},\, r_i\in \mathbb{Z}_{\geq 0}, \sum r_i=s\Big\},
\quad
\forall\, s\in \mathbb{N},
\end{gather*}
and the subspace of grade zero $\mathbf{U}(\mathfrak{a}[t])[0]=\mathbf{U}(\mathfrak{a})$.

{\bf 2.2.}~Let $\mathfrak{g}$ be a~f\/inite-dimensional complex simple Lie algebra, with Cartan subalgebra~$\mathfrak{h}$.
Let~$R$ (resp.~$R^+$) be the set of roots (resp.\
positive roots) of $\mathfrak{g}$ with respect to $\mathfrak{h}$ and $\theta \in R^+$ be the highest root in~$R$.
There is a~non-degenerate, symmetric, Weyl group invariant bilinear form $(\cdot|\cdot)$ on $\mathfrak{h}^*$, which we assume to
be normalized so that the square length of a~long root is two.
For $\alpha\in R$, $\alpha^{\vee}\in\mathfrak{h}$ denotes the corresponding co-root and we set
$d_{\alpha}=2/(\alpha|\alpha)$. For $\alpha\in R$, let $\mathfrak{g}_{\alpha}$ be the corresponding root space of
$\mathfrak{g}$ and f\/ix non-zero elements $x^{\pm}_{\alpha}\in \mathfrak{g}_{\pm\alpha}$ such that $[x^{+}_{\alpha},x^{-}_{\alpha}]=\alpha^{\vee}$.
We set $\mathfrak{n}^{\pm}=\oplus_{\alpha\in R^{+}}  \mathfrak{g}_{\pm\alpha}$.

Let $P^{+}$ be the set of dominant integral weights of $\mathfrak{g}$.
For $\lambda\in P^+$, $V(\lambda)$ be the corresponding f\/inite-dimensional irreducible $\mathfrak{g}$-module generated~by
an element $v_{\lambda}$ with the following def\/ining relations:
\begin{gather*}
x^{+}_{\alpha}v_\lambda=0,
\qquad
h v_\lambda = \langle \lambda, h \rangle v_\lambda,
\qquad
(x^{-}_{\alpha})^{\langle \lambda, \alpha^{\vee} \rangle +1}  v_\lambda=0,
\qquad
\text{for all}
\quad
\alpha\in R^+,
\quad
h\in\mathfrak{h}.
\end{gather*}

{\bf 2.3.}~A graded $\mathfrak{g}[t]$-module is a~$\mathbb{Z}$-graded vector space
\begin{gather*}
V=\bigoplus_{r\in\mathbb{Z}} V[r] \qquad
\text{such that}
\quad
(x\otimes t^s) V[r]\subset V[r+s],
\quad
x\in\mathfrak{g},
\quad
r\in\mathbb{Z},
\quad
s\in\mathbb{Z}_{\geq0}.
\end{gather*}
For $\mu\in\mathfrak{h}^*$, an element~$v$ of a~graded $\mathfrak{g}[t]$-module~$V$ is said to be of weight~$\mu$, if
$(h\otimes 1)v=\langle \mu, h\rangle v$ for all $h\in \mathfrak{h}$.
We def\/ine a~morphism between two graded $\mathfrak{g}[t]$-modules as a~degree zero morphism of
$\mathfrak{g}[t]$-modules.
For $r\in\mathbb{Z}$, let $\tau_r$ be the grade shift operator: if~$V$ is a~graded $\mathfrak{g}[t]$-module then~$\tau_r
V$ is the graded $\mathfrak{g}[t]$-module with the graded pieces shifted uniformly by~$r$ and the action of~$\mathfrak{g}[t]$ remains unchanged.
For any graded $\mathfrak{g}[t]$-module~$V$ and a~subset~$S$ of~$V$, $\langle S\rangle$ denotes the submodule of~$V$ generated by~$S$.
For $\lambda\in P^+$, $\ev_0 V(\lambda)$ be the irreducible graded $\mathfrak{g}[t]$-module such
that $\ev_0 V(\lambda)[0]\cong_{\mathfrak{g}} V(\lambda)$ and $\ev_0 V(\lambda)[r]=0$
$\forall\, r\in \mathbb{N}$.
In particular, $\mathfrak{g}[t]_{+}(\ev_0 V(\lambda))=0$.

{\bf 2.4.}~For $r,s\in\mathbb{Z}_{\geq0}$, we denote
\begin{gather*}
\mathbf{S}(r,s)=\bigg\{(b_p)_{p\geq0}: b_p\in\mathbb{Z}_{\geq0},
\;
\sum\limits_{p\geq 0}b_p =r,
\;
\sum\limits_{p\geq0}pb_p=s\bigg\}.
\end{gather*}
For $\alpha\in R^{+}$ and $r,s\in\mathbb{Z}_{\geq0}$, we def\/ine an element
$\mathbf{x}^{-}_{\alpha}(r,s)\in\mathbf{U}(\mathfrak{g}[t])[s]$~by
\begin{gather*}
%\label{prelim}
\mathbf{x}^{-}_{\alpha}(r,s)=\sum\limits_{(b_p)\in\mathbf{S}(r,s)} (x^{-}_{\alpha} \otimes 1)^{(b_0)}(x^{-}_{\alpha}
\otimes t)^{(b_1)}\cdots(x^{-}_{\alpha} \otimes t^s)^{(b_s)},
\end{gather*}
where for any non-negative integer~$b$ and any $x\in\mathfrak{g}[t]$, we understand
$x^{(b)}=x^b/b!$.

The following was proved in~\cite{G} (see also~\cite[Lemma 2.3]{CV}).
\begin{lemma}\label{garland}
Given $s\in\mathbb{N}$, $r\in\mathbb{Z}_{\geq0}$ and $\alpha\in R^{+}$, we have
\begin{gather*}
(x^{+}_{\alpha} \otimes t)^{(s)}(x^{-}_{\alpha} \otimes
1)^{(s+r)}-(-1)^s\mathbf{x}^{-}_{\alpha}(r,s)\in\mathbf{U}(\mathfrak{g}[t])\mathfrak{n}^{+}[t]\bigoplus
\mathbf{U}(\mathfrak{n}^{-}[t])\mathfrak{h}[t]_{+}.
\end{gather*}
\end{lemma}

\section{Weyl, Demazure modules and fusion product}\label{section3}

In this section, we recall the def\/initions of local Weyl modules, level one Demazure modules and fusion products.

\subsection{Weyl module}\label{section3.1}

The def\/inition of the local Weyl module was given originally in~\cite{CP}, later in~\cite{CFK} and~\cite{FL}.
\begin{definition}
Given $\lambda \in P^+$, the local Weyl module $W(\lambda)$ is the cyclic $\mathfrak{g}[t]$-module generated by an
element $w_\lambda$, with following def\/ining relations:
\begin{gather}
\mathfrak{n}^{+}[t]  w_\lambda=0,
\qquad
(h \otimes t^s)  w_\lambda = \langle \lambda, h \rangle \delta_{s,0} w_\lambda=0,
\qquad
s\geq0,
\qquad
h\in\mathfrak{h},
%\label{w1}
\nonumber
\\
(x^{-}_\alpha \otimes 1)^{\langle \lambda,
\alpha^{\vee} \rangle + 1}w_\lambda=0,
\qquad
\alpha \in R^{+}.
\label{w2}
\end{gather}
\end{definition}
\noindent We note that the relation~\eqref{w2} implies
\begin{gather}
\label{w2'}
\big(x^{-}_\alpha \otimes t^{\langle \lambda,
\alpha^{\vee} \rangle}\big)  w_\lambda=0,
\qquad
\alpha \in R^{+},
\end{gather}
which is easy to see from Lemma~\ref{garland}.
We set the grade of $w_\lambda$ to be zero; then $W(\lambda)$ becomes a~$\mathbb{Z}_{\geq 0}$-graded module with
\begin{gather*}
W(\lambda)[0] \cong_{\mathfrak{g}} V(\lambda).
\end{gather*}
Moreover, $\textup{ev}_0 V(\lambda)$ is the unique graded irreducible quotient of $W(\lambda)$.

We now specialize to the case $\lambda\in \mathbb{N}\theta$, and obtain some further useful relations that hold in~$W(\lambda)$.
\begin{lemma}
\label{g}
Let $k\in\mathbb{N}$.
The following relations hold in the local Weyl module $W((k+1)\theta)$:
\begin{enumerate}\itemsep=0pt
\item[$1)$] $(x^{-}_{\theta}\otimes 1)^{2k+1}\big(x^{-}_{\theta}\otimes t^{2k+1-i}\big)w_{(k+1)\theta}=0,
\qquad
\forall\,
0\leq i \leq k$;
\item[$2)$] $(x^{-}_{\theta}\otimes t^m)(x^{-}_{\theta} \otimes t^{m+1})  w_{(k+1)\theta} \in  \big\langle \big(x^{-}_{\theta}
\otimes t^{m+2}\big) w_{(k+1)\theta} \rangle,
\qquad
\forall\,
m\geq k$.
\end{enumerate}
\end{lemma}

\begin{proof}
To prove part (1), consider $(x^{+}_{\theta}\!\otimes t^{2k{+}1{-}i}) (x^{-}_{\theta}\otimes 1)^{2k{+}3}w_{(k{+}1)\theta}$.
Since $(x^{+}_{\theta}\!\otimes t^{2k{+}1{-}i})w_{(k{+}1)\theta}$ $=0$, we get
\begin{gather*}
\big(x^{+}_{\theta}\otimes t^{2k+1-i}\big) (x^{-}_{\theta}\otimes 1)^{2k+3}w_{(k+1)\theta}
=\big[x^{+}_{\theta}\otimes t^{2k+1-i}, (x^{-}_{\theta}\otimes 1)^{2k+3}\big] w_{(k+1)\theta}
\\
\qquad{}
=\sum\limits_{j=1}^{2k+3}  (x^{-}_{\theta}\otimes 1)^{j-1}\big(\theta^{\vee}\otimes t^{2k+1-i}\big)(x^{-}_{\theta}\otimes
1)^{2k+3-j}w_{(k+1)\theta}.
\end{gather*}
Since $(\theta^{\vee}\otimes t^{2k+1-i})w_{(k+1)\theta}=0$, we may replace $(\theta^{\vee}\otimes
t^{2k+1-i})(x^{-}_{\theta}\otimes 1)^{2k+3-j}$~by
\begin{gather*}
\big[\theta^{\vee}\otimes t^{2k+1-i}, (x^{-}_{\theta}\otimes 1)^{2k+3-j}\big]
=(-2)(2k+3-j)(x^{-}_{\theta}\otimes 1)^{2k+2-j}\big(x^{-}_{\theta}\otimes t^{2k+1-i}\big).
\end{gather*}
After simplifying, we get
\begin{gather*}
\big(x^{+}_{\theta}\otimes t^{2k+1-i}\big) (x^{-}_{\theta}\otimes 1)^{2k+3}w_{(k+1)\theta}
\\
\qquad{}
=(-1)(2k+2)(2k+3)(x^{-}_{\theta}\otimes 1)^{2k+1}\big(x^{-}_{\theta}\otimes t^{2k+1-i}\big)w_{(k+1)\theta}.
\end{gather*}
Now, using $(x^{-}_{\theta}\otimes 1)^{2k+3}w_{(k+1)\theta}=0$ in $W((k+1)\theta)$, completes the proof of part~(1).
Part (2) follows easily by putting $r=2$, $s=2m+1$ and $\alpha=\theta$ in Lemma~\ref{garland}, and using the fact that
$(x^{-}_{\theta}\otimes 1)^{2m+3}w_{(k+1)\theta}=0$, $\forall\,
m\geq k$ by~\eqref{w2}.
\end{proof}

\subsection{Level one Demazure module}\label{section3.2}

Let $\lambda\in P^{+}$ and $\alpha\in R^{+}$ with $\langle \lambda, \alpha^{\vee}\rangle > 0$.
Let $s_\alpha, m_\alpha \in \mathbb{N}$ be the unique positive integers such that
\begin{gather*}
\langle \lambda, \alpha^{\vee}\rangle = (s_\alpha-1)d_\alpha + m_\alpha,
\qquad
0<m_\alpha\leq d_\alpha.
\end{gather*}
If $\langle \lambda, \alpha^{\vee}\rangle = 0$, set $s_\alpha=0=m_\alpha$.
We take the following as a~def\/inition of the level one Demazure module.
\begin{definition}
\textup{(see~\cite[Corollary 3.5]{CV})} The level one Demazure module $D(1,\lambda)$ is the graded quotient of
$W(\lambda)$ by the submodule generated by the union of the following two sets:
\begin{gather}
\big\{(x^{-}_{\alpha} \otimes t^{s_{\alpha}})  w_{\lambda}: \alpha\in R^{+}~\textup{such that}~d_{\alpha} > 1\big\},
\label{dm1}
\\
\big\{\big(x^{-}_{\alpha} \otimes t^{s_{\alpha}-1}\big)^2  w_{\lambda}: \alpha\in R^{+}~\textup{such that}~d_\alpha =3~\textup{and}~m_\alpha=1\big\}.
\label{dm2}
\end{gather}
In particular, for $\mathfrak{g}$ simply laced, $D(1,\lambda)\cong_{\mathfrak{g}[t]} W(\lambda)$. We denote~by
$\overline{w}_\lambda$, the image of $w_\lambda$ in~$D(1,\lambda)$.
\end{definition}

The following proposition gives explicit def\/ining relations for $D(1,k\theta)$.
\begin{proposition}
\label{WvsD}
Given $k\geq 1$, the level~$1$ Demazure module $D(1,k\theta)$ is the graded $\mathfrak{g}[t]$-module generated by an
element $\overline{w}_{k\theta}$, with the following defining relations:
\begin{gather*}
\mathfrak{n}^{+}[t]\, \overline{w}_{k\theta}=0,
\qquad
(h \otimes t^s) \overline{w}_{k\theta} = \langle k\theta, h\rangle \delta_{s,0} \overline{w}_{k\theta},
\qquad
s\geq0,
\quad
h\in\mathfrak{h},
\\
(x^{-}_{\alpha}\otimes 1) \overline{w}_{k\theta}=0,
\qquad
\alpha\in R^+,
\quad
(\theta|\alpha)=0,
\\
(x^{-}_{\alpha} \otimes 1)^{kd_{\alpha}+1} \overline{w}_{k\theta}=0,
\qquad
\big(x^{-}_{\alpha} \otimes t^k\big) \overline{w}_{k\theta}=0,
\qquad
\alpha\in R^+,
\quad
(\theta|\alpha)=1,
\\
(x^{-}_{\theta} \otimes 1)^{2k+1} \overline{w}_{k\theta}=0.
\end{gather*}
\end{proposition}
\begin{proof}
Observe that, from the abstract theory of root systems $(\theta|\alpha)= 0$ or~1,
$\forall\, \alpha \in R^{+}\setminus \{\theta\}$.
This implies that $\langle k\theta, \alpha^{\vee}\rangle = 0$ or $kd_{\alpha}$, $\forall\, \alpha \in R^{+}\setminus \{\theta\}$.
Hence the relations~\eqref{dm2} do not occur in $D(1,k\theta)$ and the relations~\eqref{dm1} are
\begin{gather*}
\big(x^{-}_{\alpha} \otimes t^k\big) \overline{w}_{k\theta}=0,
\qquad
\alpha\in R^+,
\quad
\alpha~~\text{short},
\quad
(\theta|\alpha)=1.
\end{gather*}
For a~long root $\alpha\in R^+$ with $(\theta|\alpha)=1$, by~\eqref{w2'} it follows that $(x^{-}_{\alpha} \otimes t^k) \overline{w}_{k\theta}=0$.
Now the other relations are precisely the def\/ining relations of $W(k\theta)$. This proves Proposition~\ref{WvsD}.
\end{proof}
We record below a~well-known fact, for later use:
\begin{gather*}
D(1, \theta)\cong_{\mathfrak{g}} V(\theta) \oplus \mathbb{C}.
\end{gather*}
In particular,
\begin{gather}
\label{dimd}
\dim D(1, \theta)= \dim V(\theta)+1.
\end{gather}

The following is a~crucial lemma, which we use in proving Theorem~\ref{MT}.
\begin{lemma}
\label{crucial}
Let $k\geq 1$ and $0\leq i \leq k$. The following relations hold in the module $D(1,(k+1)\theta)$:
\begin{enumerate}\itemsep=0pt
\item[$1)$] $(x^{-}_{\alpha} \otimes 1)^{kd_{\alpha}+1}\big(x^{-}_{\theta} \otimes t^{2k+1-i}\big) \overline{w}_{(k+1)\theta}=0$,
$\forall\, \alpha \in R^{+}$, $(\theta|\alpha)=1$;
\item[$2)$] $\big(x^{-}_{\alpha} \otimes t^{k}\big) (x^{-}_{\theta} \otimes t^{2k+1-i}) \overline{w}_{(k+1)\theta} \in
\big\langle \big(x^{-}_{\theta} \otimes t^{2k+2-i}\big) \overline{w}_{(k+1)\theta} \big\rangle$,
$\forall\, \alpha \in R^{+}$, $(\theta|\alpha)=1$;
\item[$3)$] $\big(x^{-}_{\theta} \otimes t^{2k-i}\big) \big(x^{-}_{\theta} \otimes t^{2k+1-i}\big) \overline{w}_{(k+1)\theta}
\in
\big\langle \big(x^{-}_{\theta} \otimes t^{2k+2-i}\big) \overline{w}_{(k+1)\theta} \big\rangle$.
\end{enumerate}
\end{lemma}
\begin{proof}
Let $ \alpha \in R^{+} \textup{with} (\theta|\alpha)=1$. This implies that $\theta-\alpha$ is also a~root of
$\mathfrak{g}$ and $(\theta|\theta-\alpha)=1$. We now prove part (1).
Observe that, $(x^{-}_{\theta} \otimes t^{2k+1-i}) \overline{w}_{(k+1)\theta}$ is an element of weight $k\theta$.
Further $(x^{+}_{\alpha} \otimes 1) (x^{-}_{\theta} \otimes t^{2k+1-i}) \overline{w}_{(k+1)\theta}=0$, since
$(x^{+}_{\alpha} \otimes 1)\overline{w}_{(k+1)\theta}=0$ and $(x^{-}_{\theta-\alpha} \otimes t^{2k+1-i}) \overline{w}_{(k+1)\theta}=0$,
for all $0\leq i \leq k$.
Considering the copy of $\mathfrak{sl}_2$ spanned by $x^{+}_{\alpha} \otimes 1$, $x^{-}_{\alpha} \otimes 1$,
$\alpha^{\vee}\otimes 1$, we obtain part (1) by standard $\mathfrak{sl}_2$ arguments.
We now prove part (2).
Putting $r=2$, $s=(3k+1-i)$ and $\alpha=\theta$ in Lemma~\ref{garland}, we get
\begin{gather}
\big(x^{-}_{\theta} \otimes t^{k}\big) \big(x^{-}_{\theta} \otimes t^{2k+1-i}\big) \overline{w}_{(k+1)\theta} +
\sum\limits_{\substack{k+1\leq p\leq q \leq 2k-i\\p+q= 3k+1-i}}
\frac{1}{(2\delta_{p, q})!} \big(x^{-}_{\theta} \otimes t^{p}\big) \big(x^{-}_{\theta} \otimes t^{q}\big)
\overline{w}_{(k+1)\theta}
\nonumber\\
\qquad{}
\in
\big\langle\big(x^{-}_{\theta}\otimes t^{2k+2-i}\big) \overline{w}_{(k+1)\theta} \big\rangle, \label{e}
\end{gather}
since $(x^{-}_{\theta}\otimes 1)^{3k+3-i}\overline{w}_{(k+1)\theta}=0$, $\forall\, 0\leq i \leq k$.
Now we act on both sides of~\eqref{e} by $x^{+}_{\theta-\alpha}$ and use the relation $(x^{-}_{\alpha} \otimes t^r) \overline{w}_{(k+1)\theta}=0$,
for all $r\geq(k+1)$, which gives part~(2).
Part~(3) is immediate from the part~(2) of Lemma~\ref{g}.
\end{proof}

\subsection{Fusion product}\label{section3.3}

In this subsection, we recall the def\/inition of the fusion product of f\/inite-dimensional graded cyclic
$\mathfrak{g}[t]$-modules given in~\cite{FL} and give some elementary properties.

For a~cyclic $\mathfrak{g}[t]$-module~$V$ generated by~$v$, we def\/ine a~f\/iltration $F^{r}V$, $r\in\mathbb{Z}_{\geq 0}$~by
\begin{gather*}
F^{r}V=\sum\limits_{0\leq s \leq r} \mathbf{U}(\mathfrak{g}[t])[s] v.
\end{gather*}
We say $F^{-1}V$ is the zero space.
The associated graded space $\gr V=\bigoplus_{r\geq 0} F^{r}V/ F^{r-1}V $ naturally becomes a~cyclic
$\mathfrak{g}[t]$-module generated by $v+F^{-1}V$, with action given by
\begin{gather*}
(x\otimes t^s)\big(w+F^{r-1}V\big):= (x\otimes t^s)w+F^{r+s-1}V,
\qquad
\forall\,
x\in \mathfrak{g},
\quad
w\in F^{r}V, \quad
r, s\in \mathbb{Z}_{\geq0}.
\end{gather*}
Observe that, $\gr V \cong V$ as $\mathfrak{g}$-modules.

The following lemma will be useful.
\begin{lemma}
\label{f1}
Let~$V$ be a~cyclic $\mathfrak{g}[t]$-module.
For $r,s\in \mathbb{Z}_{\geq0}$, the following equality holds in the quotient space $F^{r+s}V/ F^{r+s-1}V$.
\begin{gather*}
(x\otimes t^s)\big(w+F^{r-1}V\big)=\big((x\otimes (t-a_1)\cdots (t-a_s))w\big) + F^{r+s-1}V,
\end{gather*}
for all $a_1,\dots,a_s\in\mathbb{C}$,
$x\in \mathfrak{g}$, $w\in F^{r}V$.
\end{lemma}

Given a~$\mathfrak{g}[t]$-module~$V$ and $z\in\mathbb{C}$, we def\/ine an another $\mathfrak{g}[t]$-module action on~$V$
as follows:
\begin{gather*}
(x\otimes t^s)v=\big(x\otimes (t+z)^s\big)v,
\qquad
x\in\mathfrak{g},
\qquad
v\in V,
\qquad
s\in \mathbb{Z}_{\geq0}.
\end{gather*}
We denote this new module by $V^z$.

Let $V_i$ be a~f\/inite-dimensional cyclic graded $\mathfrak{g}[t]$-module generated by $v_i$, for $1\leq i\leq m$, and
let $z_1,\dots,z_m$ be distinct complex numbers.
We denote~by
\begin{gather*}
\mathbf{V}={V_1}^{z_1}\otimes\dots\otimes {V_m}^{z_m},
\end{gather*}
the corresponding tensor product of $\mathfrak{g}[t]$-modules.
It is easily checked (see~\cite[Proposition 1.4]{FL}) that $\mathbf{V}$ is a~cyclic $\mathfrak{g}[t]$-module generated
by $v_1\otimes\dots\otimes v_m$.
The associated graded space $\gr \mathbf{V}$ is called the fusion product of $V_1,\dots,V_m$ w.r.t.\ the parameters
$z_1,\dots,z_m$, and is denoted by ${V_1}^{z_1}*\cdots*{V_m}^{z_m}$.
We denote $v_1*\cdots*v_m=(v_1\otimes\cdots\otimes v_m)+F^{-1}\mathbf{V}$, a~generator of $\gr \mathbf{V}$.
For ease of notation we mostly, just write $V_1*\cdots*V_m$ for ${V_1}^{z_1}*\cdots*{V_m}^{z_m}$.
But unless explicitly stated, it is assumed that the fusion product does depend on these parameters.

The following lemma will be needed later.
\begin{lemma}
\label{f2}
Given $1\leq i\leq m$, let $V_i$ be a~finite-dimensional cyclic graded $\mathfrak{g}[t]$-module generated by $v_i$, and
$s_i\in\mathbb{Z}_{\geq0}$.
Let $x\in\mathfrak{g}$.
If $(x\otimes t^{s_i})v_i=0$, $\forall\, 1\leq i\leq m$ then $(x\otimes t^{s_1+\dots+s_m}) v_1*\dots*v_m=0$.
\end{lemma}
\begin{proof}
Let $z_1,\dots,z_m$ be distinct complex numbers and let $\mathbf{V}$ as above.
By using Lemma~\ref{f1}, we get the following equality in $\gr \mathbf{V}$,
\begin{gather*}
\big(x\otimes t^{s_1+\cdots+s_m}\big) \big((v_1\otimes\cdots\otimes v_m)+F^{-1}\mathbf{V}\big)
\\
\qquad{}
=\big(\big(x\otimes (t-z_1)^{s_1}\cdots(t-z_m)^{s_m}\big)v_1\otimes\dots\otimes v_m\big) + F^{s_1+\cdots+s_m-1}\mathbf{V}.
\end{gather*}
Now the proof follows by the def\/inition of the fusion product.
\end{proof}

\section{Proof of the main theorem}\label{section4}

In this section, we prove the existence of maps $\phi^{+}$ and $\phi^{-}$ from Theorem~\ref{MT} and then prove our main
theorem (Theorem~\ref{MT}).

{\bf 4.1.}~Given $k\geq 1$ and $0\leq i \leq k $, we denote~by
\begin{gather*}
\mathbf{V}_{i,k} = D(1, k\theta)/\big\langle \big(x^{-}_\theta \otimes t^{2k-i}\big) \overline{w}_{k\theta} \big\rangle,
\end{gather*}
and let $\overline{v}_{i,k}$ be the image of $\overline{w}_{k\theta}$ in $\mathbf{V}_{i,k}$.

Using Proposition~\ref{WvsD}, $\mathbf{V}_{i,k}$ is the cyclic graded $\mathfrak{g}[t]$-module generated by the element
$\overline{v}_{i,k}$, with the following def\/ining relations:
\begin{gather}
(x^{+}_{\alpha}\otimes t^s)\overline{v}_{i,k}=0,
\qquad
s\geq 0,\quad \alpha\in R^+,
\label{r1}
\\
(h \otimes t^s) \overline{v}_{i,k} = \langle k\theta, h\rangle \delta_{s,0} \overline{v}_{i,k},
\qquad
s\geq0,
\quad
h\in\mathfrak{h},
\label{r2}
\\
(x^{-}_{\alpha}\otimes 1) \overline{v}_{i,k}=0,
\qquad
\alpha\in R^+,
\quad
(\theta|\alpha)=0,
\label{r3}
\\
(x^{-}_{\alpha}\otimes 1)^{kd_{\alpha}+1} \overline{v}_{i,k}=0,
\qquad
\big(x^{-}_{\alpha} \otimes t^k\big) \overline{v}_{i,k}=0,
\qquad
\alpha\in R^+,
\quad
(\theta|\alpha)=1,
\label{r4}
\\
(x^{-}_{\theta}\otimes 1)^{2k+1} \overline{v}_{i,k}=0,
\qquad
\big(x^{-}_{\theta} \otimes t^{2k-i}\big) \overline{v}_{i,k}=0.
\label{r5}
\end{gather}
The existence of $\phi^{+}$ is trivial, which we record below.
\begin{proposition}
The map $\phi^{+} : \mathbf{V}_{i,k+1} \rightarrow \mathbf{V}_{i+1,k+1}$ which takes $\overline{v}_{i,k+1}
\rightarrow \overline{v}_{i+1,k+1}$ is a~surjective morphism of $\mathfrak{g}[t]$-modules with $\ker \phi^{+} = \big\langle
\big({x^{-}_\theta} \otimes t^{2k+1-i}\big) \overline{v}_{i,k+1}\big\rangle$.
\end{proposition}

Now we prove the existence of $\phi^{-}$ in the following proposition.
\begin{proposition}
There exist a~surjective morphism of $\mathfrak{g}[t]$-modules $\phi^{-} : \tau_{2k+1-i} \mathbf{V}_{i,k}
\rightarrow \ker \phi^{+}$,
such that $\phi^{-}(\overline{v}_{i,k}) = \big(x^{-}_\theta \otimes t^{2k+1-i}\big) \overline{v}_{i,k+1}$.
\end{proposition}

\begin{proof}
We only need to show that $\phi^{-}(\overline{v}_{i,k})$ satisf\/ies the def\/ining relations of $\mathbf{V}_{i,k}$. We
start with the relation~\eqref{r1}.
First, for $\alpha=\theta$ it is clear.
Let $\alpha \in R^{+}\setminus\{\theta\}$; if $(\theta|\alpha)=0$ then also it is clear.
If $(\theta|\alpha)=1$ then $(\theta-\alpha)\in R^{+}\setminus\{\theta\}$ and $(\theta|\theta-\alpha)=1$, now it is
clear from the relations $(x^{-}_{\theta-\alpha}\otimes t^r)\overline{v}_{i,k+1}=0$ for all $r\geq (k+1)$ in
$\mathbf{V}_{i,k+1}$.
The relations~\eqref{r2},~\eqref{r3} are trivially satisf\/ied by $\phi^{-}(\overline{v}_{i,k})$.
Finally the last two relations~\eqref{r4},~\eqref{r5} are also satisf\/ied by $\phi^{-}(\overline{v}_{i,k})$; in fact
these are exactly the statements of Lemmas~\ref{g} and~\ref{crucial}.
\end{proof}

{\bf 4.2.}~The existence of the surjective maps $\phi^{+}$ and $\phi^{-}$, give the following:
\begin{gather}
\label{diml}
\dim \mathbf{V}_{i,k+1} \leq \dim \mathbf{V}_{i,k} + \dim \mathbf{V}_{i+1,k+1}.
\end{gather}

The following proposition helps in proving the reverse inequality.
\begin{proposition}
The map $\psi :\mathbf{V}_{i,k+1}\rightarrow D(1,\theta)^{* (k+1-i)} * \textup{ev}_0 V(\theta)^{*i}$ such that
$\psi(\overline{v}_{i,k+1}) = \overline{w}_{\theta}^{* (k+1-i)} * v_{\theta}^{*i}$ is well-defined and surjective
morphism of $\mathfrak{g}[t]$-modules.
In particular,
\begin{gather}
\label{dimg}
\dim \mathbf{V}_{i,k+1}\geq (\dim D(1, \theta))^{k+1-i} (\dim V(\theta))^{i}.
\end{gather}
\end{proposition}
\begin{proof}
We only need to show that $\psi(\overline{v}_{i,k+1})$ satisf\/ies the def\/ining relations of $\mathbf{V}_{i,k+1}$. But
they follow easily from the following relations:
\begin{gather*}
((h \otimes 1)-\langle (k+1)\theta, h\rangle)\big(\overline{w}_{\theta}^{\otimes (k+1-i)} \otimes v_{\theta}^{\otimes i}\big)=0,
\qquad
\forall\, h\in\mathfrak{h},
\\
(x^{-}_{\alpha}\otimes 1)^{\langle (k+1)\theta,
\qquad
\alpha^{\vee} \rangle + 1}\big(\overline{w}_{\theta}^{\otimes (k+1-i)}\otimes v_{\theta}^{\otimes i}\big)=0,
\qquad
\forall\, \alpha\in R^{+}
\end{gather*}
(which holds in $D(1,\theta)^{\otimes (k+1-i)} \otimes \textup{ev}_0 V(\theta)^{\otimes i}$) and $(h \otimes
t^s)\psi(\overline{v}_{i,k+1})=0, \forall\, s\geq 1$, $h\in\mathfrak{h}$
(which holds in $D(1,\theta)^{* (k+1-i)}*\textup{ev}_0 V(\theta)^{*i}$).
Further from Lemma~\ref{f2}, by using the relations
\begin{gather*}
(x^{+}_{\alpha}\otimes t^s)\overline{w}_{\theta}=0=(x^{+}_{\alpha}\otimes t^s)v_{\theta},
\qquad
\forall\, s\geq 0, \quad \alpha\in R^{+},
\\
(x^{-}_{\alpha}\otimes t)\overline{w}_{\theta}=\big(x^{-}_{\theta}\otimes
t^{2}\big)\overline{w}_{\theta}=0=(x^{-}_{\theta}\otimes t)v_{\theta}=(x^{-}_{\alpha}\otimes t)v_{\theta},
\qquad
\forall\, \alpha \in R^{+}\setminus\{\theta\},
\end{gather*}
which holds in $D(1,\theta)$ and $\textup{ev}_0 V(\theta)$.
\end{proof}

We record below a~result from~\cite{F} and use this in proving our main theorem.
\begin{proposition}\textup{\cite[Corollary 2]{F}}
\label{F}
Given $k\geq 1$, the following is an isomorphism of $\mathfrak{g}[t]$-modules,
\begin{gather*}
\textup{ev}_0 V(\theta)^{* k} \cong D(1, k\theta)/\big\langle \big({x^{-}_\theta} \otimes t^{k}\big) \overline{w}_{k\theta}\big\rangle.
\end{gather*}
\end{proposition}

{\bf 4.3.}~We now prove Theorem~\ref{MT}, proceeding by induction on~$k$.
First, for $k=1$, we prove Theorem~\ref{MT} for $0\leq i \leq 1$.
Let $i=1$, observe that $\mathbf{V}_{1,1} \cong_{\mathfrak{g}[t]} \ev_{0} V(\theta)$.
Using Proposition~\ref{F},~\eqref{dimd},~\eqref{diml} and~\eqref{dimg} this case follows.
Let $i=0$, now observe that $\mathbf{V}_{0,1} \cong_{\mathfrak{g}[t]} D(1,\theta)$.
Using part (2) of Theorem~\ref{MT} for $i=1$ and $k=1$,~\eqref{dimd},~\eqref{diml} and~\eqref{dimg} this case also
follows.
Now let $k\geq2$, and assume Theorem~\ref{MT} holds for $(k-1)$.
We prove the assertion for~$k$, proceeding by induction on~$i$.
For $i=k$, it follows from Proposition~\ref{F},~\eqref{dimd},~\eqref{diml} and~\eqref{dimg}.
Now let $i\leq (k-1)$, and assume Theorem~\ref{MT} holds for $(i+1)$.
We now prove for~$i$.
Using part (2) of Theorem~\ref{MT}, for $(i+1)$ and~$k$, also for~$i$ and $(k-1)$, and~\eqref{diml}, we get
\begin{gather*}
\dim \mathbf{V}_{i,k+1} \leq (\dim D(1, \theta))^{k-i} (\dim V(\theta))^{i+1} + (\dim D(1, \theta))^{k-i}
(\dim V(\theta))^{i}.
\end{gather*}
Together with~\eqref{dimd}, we see
\begin{gather*}
\dim \mathbf{V}_{i,k+1} \leq (\dim D(1,\theta))^{k+1-i} (\dim V(\theta))^{i}.
\end{gather*}
Now the proof of Theorem~\ref{MT} in this case follows by~\eqref{dimg}.
This completes the proof of Theorem~\ref{MT}.
\begin{thebibliography}{99}
\bibitem[1]{CFK}
Chari V., Fourier G., Khandai T., A categorical approach to {W}eyl modules,
  \href{http://dx.doi.org/10.1007/s00031-010-9090-9}{\textit{Transform. Groups}} \textbf{15} (2010), 517--549, \href{http://arxiv.org/abs/0906.2014}{arXiv:0906.2014}.

\bibitem[2]{CP}
Chari V., Pressley A., Weyl modules for classical and quantum af\/fine algebras,
  \href{http://dx.doi.org/10.1090/S1088-4165-01-00115-7}{\textit{Represent. Theory}} \textbf{5} (2001), 191--223,
  \href{http://arxiv.org/abs/math.QA/0004174}{math.QA/0004174}.

\bibitem[4]{CV}
Chari V., Venkatesh R., Demazure modules, fusion products, and $Q$-systems,
  \href{http://dx.doi.org/10.1007/s00220-014-2175-x}{\textit{Comm. Math. Phys.}}, {t}o appear, \href{http://arxiv.org/abs/1305.2523}{arXiv:1305.2523}.

\bibitem[5]{FL}
Feigin B., Loktev S., On generalized {K}ostka polynomials and the quantum
  {V}erlinde rule, in Dif\/ferential Topology, Inf\/inite-Dimensional {L}ie
  Algebras, and Applications, \textit{Amer. Math. Soc. Transl. Ser.~2}, Vol.~194, Amer. Math. Soc., Providence, RI, 1999, 61--79,
  \href{http://arxiv.org/abs/math.QA/9812093}{math.QA/9812093}.

\bibitem[6]{F}
Feigin E., The {PBW} f\/iltration, {D}emazure modules and toroidal current
  algebras, \href{http://dx.doi.org/10.3842/SIGMA.2008.070}{\textit{SIGMA}} \textbf{4} (2008), 070, 21~pages,
  \href{http://arxiv.org/abs/0806.4851}{arXiv:0806.4851}.

\bibitem[8]{G}
Garland H., The arithmetic theory of loop algebras, \href{http://dx.doi.org/10.1016/0021-8693(78)90294-6}{\textit{J.~Algebra}}
  \textbf{53} (1978), 480--551.

\end{thebibliography}
\end{document}
